% Source: https://artofproblemsolving.com/community/c3h2951263_number_theory
% Stable identity: aops-topic-2951263-q01
Take an optimal set $s_1<\dots<s_k$.

(1) $1$ is not used: removing $1$ and adding it to the largest element increases the product.

(2) There is at most one gap of $2$ and no larger gap. If there were two gaps $s_{i+1}-s_i\ge2$ and $s_{j+1}-s_j\ge2$ with $i<j$, replace $s_i,s_{j+1}$ by $s_i+1,s_{j+1}-1$. The elements stay distinct and the sum is unchanged, while the product grows because $(s_i+1)(s_{j+1}-1)-s_is_{j+1}=s_{j+1}-s_i-1>0$. A gap of at least $3$ is treated the same way within that gap.

So the set is a block of consecutive integers $\ge2$ with possibly one interior number removed. Now check by size. $k\ge7$ is impossible, since $2+\dots+8=35>30$. $k=6$: $\{2,\dots,8\}$ minus $5$ gives $\{2,3,4,6,7,8\}$ with product $8064$. $k=5$: $\{4,\dots,8\}$ gives $6720$. $k=4$: $\{6,7,8,9\}$ gives $3024$. $k\le3$: at most $10^3$ by AM--GM.

The maximum is $\boxed{8064}$.
